\section{Abstract model classes, characterization, and definability}

\begin{annexitem}{annex:p03-lindstrom-class-bridge}{thm:p03-lindstrom-class-bridge}{Model classes of the regular syntax assignment}
\begin{proof}
Isomorphism invariance gives closure of every model class. The model class
of \(\phi\land\psi\) is the intersection of their model classes, and
that of \(\neg\phi\) is the complement among structures of the same
vocabulary. A vocabulary permutation carries interpretations and sentences
together by the renaming hypothesis, so it carries the associated class to
the renamed class. If a vocabulary is enlarged without changing a sentence,
its truth depends only on its original reduct; its class in the larger
vocabulary consists of all free expansions of its old models. These prove
all five class operations, with the source's satisfaction-pair orientation.
The extra conditions are genuinely separate. In \(L_{\infty\omega}\)
the sentence \(\bigwedge_{n<\omega}\forall xP_n(x)\) has no finite
occurrence vocabulary: on a one-point domain, changing one \(P_n\)
outside any proposed finite support changes its truth while leaving the
proposed reduct unchanged. Yet the five class operations follow by the
ordinary infinitary syntax/satisfaction induction.

For relativization take the explicitly constructed assignment whose
sentences are finite Boolean combinations of first-order sentences and
one sentence asserting finiteness of the whole carrier. It has the five
class operations, but cannot relativize that finiteness sentence to a
unary \(U\). If it could, then on infinite carriers replacement of the
carrier-finiteness sentence by falsity in its Boolean syntax would give
a first-order sentence \(\psi\) defining finiteness of \(U\).
The first-order set \(\{\psi\}\cup\{\text{at least }n\text{ members
of }U:n<\omega\}\) is finitely satisfiable on infinite carriers with
large finite \(U\). Compactness gives a model whose \(U\) and carrier
are infinite, contradicting that defining property. Finally the actual
\(Q_0\) and \(Q_1\) constructions proved in
\cref{thm:p03-lindstrom-assumptions} respectively separate the five
operations from compactness and from countable LS. Thus no converse to
the regular-package implication is silently inferred.
\end{proof}
\end{annexitem}

\begin{annexitem}{annex:p03-lindstrom-generalized-class}{thm:p03-lindstrom-generalized-class}{The actual quantifier semantics and isomorphism induction}
\begin{proof}
Define the quantifier by the displayed class-membership clause. For the
atomic matrix \(R(x,y)\), the relation selected in that clause is exactly
the relation of the input structure, proving the defining sentence assertion.
If \(h:M\cong N\) is a bijection transporting all base interpretations,
induction on terms and formulas proves ordinary first-order invariance.
In the new quantifier case the induction hypothesis transports the selected
binary relation to its image under \(h\). Those two selected structures
are isomorphic, so membership in \(K\) agrees. The quantifier formula
therefore is invariant as well; nesting causes no additional case beyond
this same induction. Conversely satisfaction invariance immediately makes
the models of a sentence isomorphism closed. This is a genuine interpretation
of the new quantifier, but adds no axiom or construction establishing
compactness or an LS property.
\end{proof}
\end{annexitem}

\begin{annexitem}{annex:p03-lindstrom-ordinal-code}{thm:p03-lindstrom-ordinal-code}{Transfinite rank, uniqueness, and the necessary order guard}
\begin{proof}
Induct on \(\alpha\). At zero the formula asserts that there are no
predecessors, precisely order type zero. Assume the predecessor-segment
claim for every \(\beta<\alpha\). If \(\eta_\alpha(a)\) holds,
every predecessor has a smaller ordinal predecessor type, and every smaller
ordinal type occurs. A point of a fixed predecessor type is unique: if two
distinct points of type \(\beta\) were comparable, the lower would be a
point in an initial segment of type \(\beta\) whose own predecessor
segment had type \(\beta\), impossible in that ordinal. Thus choose
the unique point \(b_\beta<a\) of type \(\beta\). The map
\(\beta\mapsto b_\beta\) is onto the predecessor segment.
It preserves order: for \(\beta<\gamma\), placing \(b_\gamma\)
below \(b_\beta\) would put type \(\gamma\) inside the ordinal
\(\beta\), where all predecessor types are strictly below \(\beta\).
Linearity and distinctness therefore force \(b_\beta<b_\gamma\).
The map is an order isomorphism from \(\alpha\) to that segment.
Conversely if the segment has type \(\alpha\), its point at each
\(\beta<\alpha\) has predecessor segment of type \(\beta\).
By induction these are exactly the ranks used by the recursive formula,
proving both its coverage and existence conjuncts.

For \(\eta^0_\alpha\), the same uniqueness and order arguments apply
to the entire domain rather than a predecessor segment. Coverage and
existence give an order isomorphism with \(\alpha\); the converse
follows by evaluating each ordinal's predecessor segment. At \(\alpha=0\)
the coverage condition is false on every nonempty domain, as required.
Without linearity, take two points and the empty relation \(<\).
Both points satisfy \(\eta_0\), so every point has rank zero and at
least one does; hence \(\eta^0_1\) holds. Its domain has two points,
not one. This is why the actual expansion sentence includes
\(\operatorname{LO}(<)\).
\end{proof}
\end{annexitem}

\begin{annexitem}{annex:p03-lindstrom-mcgee-code}{thm:p03-lindstrom-mcgee-code}{An enumerating block and the full signed relation diagram}
\begin{proof}
Fix a bijection \(f_A:\lambda\to A\), and distinct variables
\(x_\alpha\), \(\alpha<\lambda\). For each pair define the signed
atom \(\rho_{\alpha,\beta}(x,y)\) to be \(R(x,y)\) if
\(A\models R(f_A(\alpha),f_A(\beta))\), and \(\neg R(x,y)\)
otherwise. Both arguments use the same enumeration of \(A\).
Form the actual sentence
\[
 \theta_A:=\exists(x_\alpha)_{\alpha<\lambda}
 \left[
 \bigwedge_{\alpha<\beta<\lambda}x_\alpha\ne x_\beta
 \land\forall y\bigvee_{\alpha<\lambda}y=x_\alpha
 \land\bigwedge_{\alpha,\beta<\lambda}
                         \rho_{\alpha,\beta}(x_\alpha,x_\beta)
 \right].
\]
Its block length and all conjunction sizes are at most \(\lambda\),
so it belongs to \(L_{\lambda^+\lambda^+}\). In a model of the
sentence the witnesses are distinct and cover its domain; the map
\(f_A(\alpha)\mapsto x_\alpha\) is a bijection. The entire signed
diagram says exactly that this map preserves and reflects \(R\),
so it is an isomorphism. Conversely an isomorphism supplies such a
witnessing enumeration and every signed instance. This proves the
characteristic assertion, not just one direction of a diagram test.

There are \(2^\lambda\) binary relations on domain \(\lambda\),
because \(|\lambda^2|=\lambda\) for infinite \(\lambda\) under AC.
Select the set of those codes belonging to \(K\), and take the
disjunction of their \(\theta_A\). Each model satisfying it is
isomorphic to a selected code and hence belongs to \(K_\lambda\).
Every member of \(K_\lambda\) has such a code by transporting its
relation along a bijection to \(\lambda\), hence satisfies that
disjunction. Its index set has size at most \(2^\lambda\), giving
the stated larger connective bound. If no code is selected use falsity.
For a finite nonzero cardinal \(n\) the same finite diagram and finite
disjunction are first-order sentences; for zero the model class is empty
under the nonempty-domain convention. Thus no cardinal case needs a
proper-class formula.
\end{proof}
\end{annexitem}

\begin{annexitem}{annex:p03-lindstrom-delta-code}{thm:p03-lindstrom-delta-code}{Source ordinal positions and a global complement pair}
\begin{proof}
Fix infinite \(\lambda\). For each code \(A\) on domain
\(\lambda\), use the source signed atoms just defined and put
\[
 \Phi_A:=\forall x\forall y\bigwedge_{\alpha,\beta<\lambda}
 \left((\eta_\alpha(x)\land\eta_\beta(y))
                                      \to\rho_{\alpha,\beta}(x,y)\right),
 \qquad
 \Theta_K:=\bigvee\{\Phi_A:A\in K,\ \operatorname{dom}(A)=\lambda\}.
\]
These have finite quantifier blocks. If an expansion satisfies
\(\operatorname{LO}(<)\land\eta^0_\lambda\), the ordinal-code theorem
gives a unique point at every \(\alpha<\lambda\), exhausting the
domain. Consequently \(\Phi_A\) makes the bijection from these
positions to the coded elements an isomorphism of the binary relations.
Conversely pull the ordinal order back along any isomorphism with a
code \(A\); all position formulas and all signed atoms then hold.
Thus the projection of
\[
 \operatorname{LO}(<)\land\eta^0_\lambda\land\Theta_K
\]
onto the base vocabulary is exactly \(K_\lambda\). This is the
positive source construction with its implicit linear-order regime
made explicit.

Replacing the selected code set by its complement among all codes on
\(\lambda\) gives \(\Theta_{\neg K}\). On models of size
\(\lambda\), projection of
\(\operatorname{LO}(<)\land\eta^0_\lambda\land\Theta_{\neg K}\)
is precisely nonmembership. The source's universal-order formulation is
also exact on that cardinal: every domain of size \(\lambda\) admits
an ordinal order of type \(\lambda\), and any such order enumerates
one of its relation codes. Therefore membership holds if and only if
every order of that type makes \(\Theta_K\) true. Outside the cardinal
that universal assertion can be vacuous, so it alone is not the global
complement pair promised in the statement.

Here is the missing off-cardinality construction. Add also a unary
predicate \(P\), and let the negative expansion sentence be the
disjunction of three branches:
\[
\begin{split}
 &\bigvee_{0<\mu<\lambda,\ \mu\text{ a cardinal}}
                    (\operatorname{LO}(<)\land\eta^0_\mu),\\
 &\operatorname{LO}(<)^P\land(\eta^0_{\lambda^+})^P,\\
 &\operatorname{LO}(<)\land\eta^0_\lambda\land\Theta_{\neg K}.
\end{split}
\]
Relativization here is syntactic: restrict universal quantifiers to \(P\)
and existential quantifiers by conjunction with \(P\); it commutes
with negation and the set conjunctions. Induction on the well-founded
formula tree proves that this evaluates the original formula in the
nonempty induced \(P\)-domain. The second branch itself forces that
domain nonempty and of size \(\lambda^+\). The first branch has an
expansion exactly when the whole domain has a nonzero cardinal below
\(\lambda\). The second has an expansion exactly when its cardinal
is at least \(\lambda^+\): choose such a subset and its ordinal order
under AC; conversely its order code supplies an injection of
\(\lambda^+\) into the domain. The last branch handles precisely
the size-\(\lambda\) models outside \(K\). Since all domains are
well orderable in this metatheory, the three projections exhaust the
complement of \(K_\lambda\), disjoint from its positive projection.
The positive expansion may ignore \(P\). Both defining formulas are
set-sized members of \(L_{\infty\omega}\); hence the two projections
prove \(\Delta\)-definability in the stated global sense.

For a finite relational vocabulary, at an arity-\(k\) symbol replace
the two variables in \(\Phi_A\) by a block of \(k\) variables and
use the conjunction of their \(\eta_{\alpha_j}\) position tests.
Choose the positive or negative \(R\)-atom by that code's entire
arity-\(k\) relation table. The unique position bijection then preserves
and reflects that symbol, exactly as in the binary proof. Take the
finite conjunction over symbols; a nullary symbol has the empty argument
block and its one signed truth atom. This supplies every relation of
the same isomorphism, proving the arity bridge. For a finite cardinal
the position and code families are finite and the smaller/larger branches
are the ordinary finite cardinality sentences, giving the same conclusion.
\end{proof}
\end{annexitem}

\begin{annexitem}{annex:p03-lindstrom-absolute-truth}{thm:p03-lindstrom-absolute-truth}{Set recursion tables over finite assignments}
\begin{proof}
Encode a formula by its well-founded set of nodes, connective labels,
immediate-subformula sets, and finite variable lists. A witness to syntactic
admission includes an ordinal-valued rank function strictly decreasing
along immediate-subformula edges and finite free-variable supports at each
node. The labels, arities, variable bindings, and support equations are
checked by bounded quantification over these supplied sets; a conjunction's
support is the union of its children's supports, which must be finite.
Atomic nodes have their indicated finite relational arity. Negation preserves
support; quantification removes its bound variable. Existence of these
sets and rank/support functions is a \(\Sigma_1\) definition. Conversely
a well-founded admitted syntax tree supplies them by well-founded recursion,
so this is a syntax definition, not an extra restriction on a completed tree.
No cardinal bound on the size of a conjunction is tested in this
\(L_{\infty\omega}\) definition.

Given a structure with carrier \(M\), form the sets of finite sequences
from \(M\). They are constructed recursively from the empty sequence
by adjoining one member of \(M\), and their union is a set. For each
syntax node take all assignments of its finite support into \(M\).
The rows consisting of a node and such an assignment form a set.
By recursion on the rank define a truth-value table on these rows.
An atomic row is determined by equality of its object values or membership
in its interpreted relation. Negation exchanges the two truth values;
a conjunction is true exactly when every indicated child row is true,
and a disjunction when at least one is true. At an existential node the
value is true exactly when some \(a\in M\) makes its matrix row true,
extending the old assignment by that variable and restricting it to the
child's support. Universals follow by negation or by the bounded universal
clause. These are quantifications over the supplied child sets, assignment
sets and \(M\), not over an external \(M^\omega\).

Well-founded induction proves the table unique: at any node the already
unique child values uniquely determine each parent value. ZF recursion
proves existence. Coding the finite-sequence pools, row set and the table
as existential witnesses, all their defining equations and truth clauses
are bounded checks. Thus existence of a correct table with root value
one is \(\Sigma_1\), and existence with root value zero is also
\(\Sigma_1\). Existence and uniqueness make these complementary,
so satisfaction is \(\Delta_1\).

In transitive universes containing the admitted code and structure, syntax
children and relation interpretations are the same sets; finite assignments
are the same finite functions. Induction on the well-founded code therefore
gives identical atomic, Boolean and quantifier values. Equivalently each
universe's correct table has the unique value just proved. This establishes
the claimed truth absoluteness, under ZF; no stronger assertion about an
unread weaker set-theory axiomatization is used.
\end{proof}
\end{annexitem}

\begin{annexitem}{annex:p03-lindstrom-delta-absoluteness}{thm:p03-lindstrom-delta-absoluteness}{Witness definitions do not certify global exhaustivity}
\begin{proof}
For the positive class, existentially choose interpretations of the finitely
or set-many added symbols on the fixed carrier and a correct truth table
for the positive sentence. The expansion interpretations are sets, and the
preceding \(\Sigma_1\) truth definition makes the combined assertion
\(\Sigma_1\). The negative sentence gives the same form for the other
projection. When these classes are complementary on the declared regime,
these are respectively a definition of membership and of its negation.

Between transitive universes admitting this same complementary pair, a
positive expansion/table witness in the smaller universe persists in the
larger, and so does a negative witness. Every base model belongs to one
projection in the smaller universe, and disjointness in the larger prevents
that truth value from changing. Hence membership agrees. But deciding
that two arbitrary projected sentences cover \emph{every} base structure
requires a global model-class assertion. It is not among the bounded syntax
and table checks just supplied. For the cardinal-dependent construction,
even the assertion that the fixed ordinal parameters denote the intended
cardinals is a separate ambient condition. Therefore the semantic witness
proof does not infer a \(\Sigma_1\) predicate for admission of all
\(\Delta\)-sentence pairs or unqualified preservation of that admission
under changes of universe. This is exactly the syntax/semantics distinction
needed for the source's stated absoluteness limitation.
\end{proof}
\end{annexitem}

\begin{annexitem}{annex:p03-lindstrom-amalgamation}{thm:p03-lindstrom-amalgamation}{First-order translations and the two full elementary diagrams}
\begin{proof}
Record explicitly the finite-rank/game bridge used by the imported
characterization. For finite relational vocabulary and a fixed finite
variable list, the bounded-rank formulas have finitely many equivalence
classes, by the rank induction supplied there. Choose their finite list
of representatives. The conjunction which signs every representative
according to its truth at a tuple characterizes that tuple's entire
bounded-rank type. If tuples agree at rank \(r+1\), an extension on
one side has a rank-\(r\) characteristic formula; the existential
closure of that formula is true there and has rank at most \(r+1\).
Agreement supplies an extension on the other side satisfying it. Thus
the enlarged tuples agree at rank \(r\). Repeating in either direction
gives a strategy for \(r+1\) rounds, ending at atomic agreement, which
is exactly a partial isomorphism in the relational vocabulary.
Conversely induct on rank and formula construction for a winning finite
strategy: atomic formulas agree at every position, Boolean combinations
retain agreement, and a witness for an existential matrix is matched by
the next strategy move. The remaining strategy and induction give
agreement on the lower-rank matrix. Negation handles the universal case.
This proves that bounded-rank agreement is equivalent to a winning
finite back-and-forth strategy; it is not inferred merely from the
existence of some partial maps.

Apply the already supplied regular-logic characterization
\cref{thm:p03-lindstrom}, whose complete coding, finite-rank and
back-and-forth construction is \cref{annex:p3.abstract.3}.
Each sentence \(\phi\) has a first-order equivalent over its finite
occurrence vocabulary. Enlarging that vocabulary preserves equivalence
by reduct invariance. For a set of sentences choose one such first-order
equivalent for each, using AC. Every finite set of chosen equivalents is
satisfiable exactly when the corresponding original finite set is.
First-order compactness gives a model of all the equivalents, hence of
all the original sentences. This proves set-sized compactness from the
stated characterization, not from countable compactness alone.

For a map, test sentences after naming each finite parameter tuple.
The characterization applies uniformly to these finite-constant expansions
by its declared abstraction clause. A first-order elementary map therefore
preserves and reflects the equivalent first-order sentence for every
\(\mathcal L^*\)-sentence with such parameters. Conversely the logic
contains first-order sentences, so \(\mathcal L^*\)-elementarity
implies first-order elementarity. This proves the semantic map bridge.

Assume now \(C\preccurlyeq A,B\) by that bridge. Give all members
of \(A\) and \(B\) parameter names, with the same names on \(C\)
and disjoint new names elsewhere, and unite their full first-order
elementary diagrams. A finite diagram fragment from \(A\) is a
conjunction \(\alpha(\bar a,\bar c)\); one from \(B\) is
\(\beta(\bar b,\bar c)\). Each can have additional finite parameter
tuples from \(C\), included in the displayed \(\bar c\). Since
\(C\) is elementary in both sides,
\(C\models\exists\bar x\alpha(\bar x,\bar c)\) and
\(C\models\exists\bar y\beta(\bar y,\bar c)\). Choose both
witness tuples in \(C\), independently for the disjoint new names.
They make that finite union true in \(C\). Compactness supplies a model
\(D\) of both entire diagrams. The diagram tests give elementary
embeddings of \(A,B\) into \(D\), agreeing on \(C\) because its
names were shared. The map bridge makes them \(\mathcal L^*\)-elementary.
No disjointness of the two images outside \(C\) is asserted.
For countable inputs and vocabulary, downward Lowenheim--Skolem supplies
a countable elementary submodel containing both images; the bridge again
preserves the desired elementarity. This is a construction for the actual
specified models, unlike a joint-consistency assertion about two theories
without a prescribed common model.
\end{proof}
\end{annexitem}

\begin{annexitem}{annex:p03-lindstrom-assumptions}{thm:p03-lindstrom-assumptions}{A countable closure carrying all required infinite witnesses}
\begin{proof}
Use the syntax generated from first-order formulas by the \(Q_0\)
quantifier, with its displayed at-least-countably-infinite semantics.
Every expression has finitely many symbols, so it has finite occurrence
vocabulary. Isomorphism, renaming and reduct invariance follow by syntax
induction; in the new clause a bijection carries its witness set bijectively
to the corresponding witness set. Boolean operations are the standard
ones. Fresh-constant abstraction has the usual existential semantics.
For unary relativization replace \(Q_0x\phi\) by
\(Q_0x(U(x)\land\phi^U)\), and use ordinary guarded first-order
quantifiers. Induction on the matrix and equality of the selected witness
sets prove truth in the nonempty induced \(U\)-structure. These give
the regular relational operations, not a promised compactness theorem.

Let \(M\) be a model of a sentence in countable vocabulary. Starting
with one element and all constant values if any, build countable sets
\(A_0\subseteq A_1\subseteq\cdots\) as follows. For every formula
\(\exists x\phi(x,\bar y)\) and tuple \(\bar a\) from \(A_n\),
if it is true in \(M\), add one actual witness. For every formula
\(Q_0x\phi(x,\bar y)\) and tuple from \(A_n\), if it is true in
\(M\), add an actual sequence of pairwise distinct witnesses. Such a
sequence exists under AC. There are countably many formulas and finite
tuples at each stage, so only countably many elements are added. Include
function values if a relational presentation has been expanded by constants
or functions; their finite arities preserve this count. Put
\(N=\bigcup_nA_n\), with the induced interpretations.

Induct on \(L(Q_0)\)-formulas to show that their truth on tuples in
\(N\) agrees with their truth in \(M\). Atomic and Boolean clauses
are the induced-structure clauses. A true existential statement in \(M\)
has its scheduled witness in \(N\); a witness in \(N\) remains true
in \(M\) by induction. For \(Q_0\), a true statement in \(M\)
has a scheduled infinite set of distinct matrix witnesses in \(N\),
true there by induction. If it is false in \(M\), its matrix witness
set is finite; induction makes the witness set in \(N\) its intersection
with \(N\), still finite. Conversely an infinite witness set in
\(N\) is one in \(M\). These prove the new quantifier clause and
complete the induction. Hence the original sentence has a countable
model, with no appeal to an unproved expanded-language Skolem theorem.

Nevertheless
\(\{\neg Q_0x(x=x)\}\cup\{\text{at least }n\text{ elements}:n<\omega\}\)
is finitely satisfiable in sufficiently large finite structures and has
no model: all finite lower bounds force infinity, which the first sentence
forbids. This disproves countable compactness. If the infinity sentence
were first-order definable, the same set with its first-order negation
would contradict first-order compactness; thus this is a proper expressive
extension. Finally \(Q_1x(x=x)\) has a model on \(\omega_1\) but none
on any countable domain by its defined semantics. The isomorphism and
five class operations for that quantifier follow from the same bijection
induction; only its asserted LS failure is used here.
\end{proof}
\end{annexitem}

\begin{annexitem}{annex:p03-lindstrom-two-point-code}{thm:p03-lindstrom-two-point-code}{Four bits, two enumerations, and incompatible diagrams}
\begin{proof}
Under \(\operatorname{LO}(<)\land\eta^0_2\), there are exactly two
positions, uniquely tested by \(\eta_0\) and \(\eta_1\). The first
two signed instances make the position-zero element related to itself
and to the other element; the two negative instances forbid both outgoing
relations from position one. These four cases are the entire relation
table, so the ordinal position map is an isomorphism with the given
structure. Conversely place the loop vertex first in any isomorphic copy;
the order expansion then satisfies exactly those four signed instances.

On labeled domain \(2\) there are sixteen binary relations. Exactly
two of their matrices are codes for this isomorphism type, corresponding
to the two choices of labeling the loop vertex. The positive projection
can use the displayed matrix alone, because the auxiliary order can be
chosen with that vertex first. The negative size-two branch uses the
disjunction of the other fourteen matrices with their ordinal tests.
Add the smaller-domain branch
\(\operatorname{LO}(<)\land\eta^0_1\) and the larger-domain branch
\(\operatorname{LO}(<)^P\land(\eta^0_3)^P\). The latter has an
expansion exactly when there are at least three elements. These branches
are set-sized, actually finite, and together project to exactly the
global complement of the positive two-point class. In particular the
relabeled copy is not incorrectly placed in the negative branch.

For the failed amalgamation premise, use vocabulary \(\{P\}\).
Let \(C=A\) have one point and empty \(P\); let \(B\) extend that
domain by a point \(b\) in \(P\), with \(P\) still empty on \(C\).
Both inclusions of \(C\) are substructure inclusions. The elementary
diagram of \(A\) contains \(\forall x\neg P(x)\), while that of
\(B\) contains \(P(c_b)\). Their union is inconsistent, so no common
target admits both elementary embeddings. The premise fails precisely
because \(C\) is not elementary in \(B\), not because compactness
or the proved elementary amalgamation construction fails.
\end{proof}
\end{annexitem}
